\section{Metodología del proyecto}

La estrategia general de este proyecto se fundamenta en la metodología \textbf{Design Science Research (DSR)}, la cual permite la creación y evaluación de artefactos tecnológicos diseñados para resolver problemas prácticos identificados en la ingeniería de software \parencite{Hamed2023}. En este contexto, el problema abordado es la falta de gobernanza técnica continua en PyMES, y el artefacto propuesto es un mecanismo automatizado de detección de deriva arquitectónica.

El enfoque del trabajo es eminentemente técnico y se divide en cuatro fases iterativas, alineadas directamente con los objetivos específicos planteados.

\subsection{Fases del Trabajo}

\begin{enumerate}
    \item \textbf{Fase de Definición y Especificación (Objetivo 1):} Se procederá a la formalización de las reglas de diseño bajo el paradigma de la arquitectura evolutiva. Utilizando el concepto de \textit{Fitness Functions} \parencite{PaulWang2019}, se traducirán principios abstractos (como la integridad de capas y la modularidad) en indicadores métricos objetivos y ejecutables. Se priorizarán métricas de acoplamiento eferente y detección de ciclos por ser síntomas críticos de erosión \parencite{Pigazzini2021}.

    \item \textbf{Fase de Implementación del Analizador Estático (Objetivo 2):} Se desarrollará un componente capaz de extraer hechos arquitectónicos directamente del código fuente Java sin requerir ejecución. Siguiendo modelos de inferencia de diseño \parencite{Klein2022}, el componente procesará archivos .java y configuraciones de construcción (Maven/Gradle). Los datos extraídos se representarán mediante \textbf{grafos de dependencia}, donde las clases y paquetes actuarán como nodos y sus relaciones como aristas, utilizando bases de datos orientadas a grafos para facilitar consultas estructurales complejas \parencite{Pigazzini2021, Klein2022}.

    \item \textbf{Fase de Construcción del Motor de Validación (Objetivo 3):} Se programará el motor lógico que contrastará el grafo de la arquitectura implementada (extraído en la fase anterior) contra el modelo de la arquitectura prescrita definido por las \textit{Fitness Functions}. El resultado principal será el cálculo automático de la \textbf{Tasa de Conformidad (CR)}, que proporcionará una visión cuantitativa de la integridad del sistema \parencite{PaulWang2019}.

    \item \textbf{Fase de Evaluación mediante Estudio de Caso (Objetivo 4):} El mecanismo se validará en un entorno controlado (real o simulado) con un esfuerzo de 192 horas de investigación. La técnica de recolección de datos consistirá en la \textbf{inyección voluntaria de derivas arquitectónicas} (violación de capas y ciclos de dependencia) para medir la sensibilidad y precisión del motor en su detección. Se utilizarán índices de referencia como el \textit{Architectural Technical Debt Index} (ATDI) para contrastar la efectividad del prototipo \parencite{SasAvgeriou}.
\end{enumerate}

\subsection{Justificación de las Decisiones Metodológicas}

La elección de \textbf{Design Science Research} se justifica por su capacidad para cerrar la brecha entre la teoría académica y la práctica industrial en PyMES, donde la documentación suele ser inexistente o estar desactualizada \parencite{Hamed2023, Anthony2024}. 

El uso de \textbf{análisis estático} basado en el código como "única fuente de verdad" es una decisión crítica para garantizar que la gobernanza sea viable en entornos ágiles, evitando la dependencia de artefactos externos que introducen errores de interpretación \parencite{Anthony2024, Klein2022}. Finalmente, la adopción de \textbf{Fitness Functions} permite integrar la validación arquitectónica directamente en el flujo de desarrollo, transformando la gobernanza técnica de una revisión manual esporádica a un proceso de monitoreo continuo y objetivo \parencite{PaulWang2019, NiessenMedium}.
